\section{Introduction}
\label{sec:intro}

Frontier language models are increasingly capable of conducting long-horizon automated research, repeatedly proposing changes, running experiments, interpreting feedback, and refining executable artifacts~\citep{MLAgentBench2024Huang, REbench2025Wijk, xu2026autolab}.
By requiring agents to optimize models, algorithms, or computing systems, these tasks can provide a measurable form of AI-for-AI and an early window into how close frontier language models are to enabling recursive self-improvement~\citep{chan2025mlebench, posttrainbench_2026, lyu2026mlsbench}.
Systematically evaluating current agents is therefore essential for understanding their research capabilities and guiding targeted improvements to both models and agent systems~\citep{REbench2025Wijk, meng2026scientistone}.


However, there remains a fundamental mismatch between the process of automated research and its evaluation: agents engage in long-horizon, closed-loop cycles of experimentation and refinement, yet existing benchmarks primarily evaluate them using a single final score, which fails to capture the underlying reasons for model behavior or provide fine-grained diagnostic information~\citep{MLAgentBench2024Huang, chan2025mlebench, REbench2025Wijk}.
Within a run, for example, the same final score may come from an effective direction identified early or from one found after extensive trial and error.
It also does not show whether proposed ideas are translated into reliable implementations or whether feedback is used effectively to retain progress and recover from failures.
Beyond the limitations of final-score-based evaluation, conventional evaluation treats each run independently, making capability appear static and obscuring whether accumulated experience improves or misleads subsequent decisions.
It also leaves unclear whether the surrounding harness helps the agent sustain effective behavior over a long research run.
These gaps make it difficult to determine how far current agents have progressed toward autonomous research and whether further improvements should target model capabilities, experience reuse, or harness design.


To address these gaps and provide a systematic evaluation framework for agents engaged in long-horizon AI R\&D, we organize the evaluation around four questions:
\begin{enumerate*}
    \item[\ding{182}] How strong are the final results produced by current agents?
    \item[\ding{183}] Where is progress gained or lost within the research loop?
    \item[\ding{184}] Can accumulated experience improve subsequent decisions?
    \item[\ding{185}] How does harness choice affect agent performance?
\end{enumerate*}


Specifically, final performance is measured directly from task scores.
To diagnose behavior within a run, we decompose the research process into three complementary capabilities: Solution Framing (C1), Execution (C2), and Feedback Control (C3).
For each capability, we design a rule-based metric computed deterministically from verifier outcomes and recorded trajectory signals rather than LLM judgments.
Beyond these within-run capabilities, we treat the ability to use accumulated experience as a meta-capability (M) and conduct controlled comparisons to measure its effect on subsequent decisions in intra- and inter-task settings.
Figure~\ref{fig:evalution_framework} summarizes the proposed process and experience views.
In addition, harness effects are examined by comparing alternative harness designs.
As a complementary analysis, we use LLM judges to examine whether the solutions produced by agents exhibit genuine methodological novelty. To provide a comprehensive evaluation, we conduct these analyses across seven frontier models and a suite of 36 long-horizon tasks. 
The full evaluation required approximately one hundred thousand U.S. dollars in model inference.

\begin{figure*}[t]
    \centering
    \includegraphics[width=\linewidth]{Figures/intro_figure_refined_v2.pdf}
    \caption{Analytical views used to interpret behavior in automated research. The process view covers Solution Framing (C1), Execution (C2), and Feedback Control (C3). The experience view uses controlled comparisons to measure how accumulated experience affects subsequent decisions in intra- and inter-task settings.}
    \label{fig:evalution_framework}
\end{figure*}

Benefiting from our evaluation design, which enables fine-grained diagnosis of model capabilities, we reach an overall assessment of the current capability stage:
\textbf{Current automated research agents operate more like engineering optimizers than fully autonomous researchers.}
Within bounded research loops, they can formulate practical directions, implement working solutions, and improve technical artifacts.
Yet their success varies across runs, genuine algorithmic innovation remains rare, and realized performance is shaped by process bottlenecks, accumulated experience, and harness design.
The evidence for this conclusion is threefold:
\begin{itemize}
    \item \textbf{Reliability separates current models more than peak performance.}
    The gap between the strongest and weakest models is $0.237$ on avg@3 but only $0.122$ on best@3.
    Several models can therefore reach competitive solutions, but show substantially different levels of consistency across repeated runs.
    These results point to headroom in inference-time selection and rollout-relative training to close the gap between observed peak and average performance.

    \item \textbf{Outcome scores conceal where research actually fails.}
    For example, GPT-5.5 and Gemini-3.1-Pro achieve similar final scores and identical Solution Framing scores, yet GPT-5.5 is substantially stronger in Execution while Gemini-3.1-Pro is stronger in Feedback Control.
    The dominant bottleneck also varies by task category: CUDA tasks show the weakest Solution Framing and Execution, whereas Model Development tasks show the strongest Execution but the weakest Feedback Control.
    More broadly, although Execution scores are high across all seven models, only three of 252 best-seed solutions qualify as novel approaches under our review protocol, revealing a clear gap between optimization performance and methodological novelty.
    Overall, a leaderboard can rank systems, but it cannot diagnose where improvement is needed.


    \item \textbf{Research performance is not fixed by the backbone model alone: experience can improve or degrade performance, while harnesses mainly affect reliability.}
    Within tasks, accumulated experience usually improves the next solution by preserving useful discoveries, but it can also carry forward misleading conclusions or anchor agents to local optima.
    Across tasks, this dual effect is strong enough to change model ordering: transferred experience raises DeepSeek-V4-Pro's avg@3 by $0.093$ but lowers Gemini-3.1-Pro's by $0.017$.
    By contrast, using their native harnesses gives GPT-5.5 and Kimi-K2.7-Code greater run-to-run stability than the shared harness without materially changing best@3 or model ordering.
    Automated harness optimization offers further headroom.
    Evaluating models as static, isolated components therefore misses both their learning dynamics and the system design required to realize their capabilities.
    
\end{itemize}




Taken together, these findings show that autonomous research capability is neither one-dimensional nor static, and that observed performance reflects an interaction among the model, its accumulated experience, and the system around it.
The report therefore begins with final performance and cost, then examines within-run process behavior, experience-driven improvement, and harness effects before discussing their implications for model training and agent-system design.








